\section{Pre-emphasis and De-emphasis in FM}

\subsection{Why FM Needs Emphasis}
\label{sec:why_fm_needs_emphasis}

In \keyterm{Frequency Modulation (FM)}, the output noise power spectral density grows as $f^2$, so the high audio frequencies (treble above about 2.2 kHz) are hit hardest by noise. Speech and music carry little power there, so the $\mathrm{SNR}$ at high frequencies is poor.

\textbf{Idea:} boost the high frequencies before modulation and cut them by the same amount after demodulation.

\begin{itemize}
    \item \keyterm{Pre-emphasis} = high-pass boost at the transmitter.
    \item \keyterm{De-emphasis} = matching low-pass cut at the receiver.
    \begin{itemize}
        \item The signal comes out unchanged.
        \item The high-frequency noise is reduced by about 13 dB.
    \end{itemize}
\end{itemize}

\subsection{Transfer Functions}
\label{sec:emphasis_transfer_functions}

The pre-emphasis filter is
\begin{equation}
    H_p(f) = 1 + j\frac{f}{f_1}
    \label{eq:pre_emphasis_filter}
\end{equation}
where:
\begin{itemize}
    \item $f_1 = \dfrac{1}{2\pi RC}$ is the corner frequency
    \item $RC = 75\ \mu\text{s}$ gives $f_1 \approx 2.1$ kHz
\end{itemize}

The de-emphasis filter is the inverse of \cref{eq:pre_emphasis_filter}, so the cascade is flat:
\begin{equation}
    H_p(f)\,H_d(f) = \left(1 + j\frac{f}{f_1}\right)\frac{1}{1 + j\frac{f}{f_1}} = 1
    \label{eq:emphasis_cascade}
\end{equation}

\begin{figure}[H]
    \centering
    \includegraphics[width=0.9\linewidth]{pre_de_emphasis_chain.png}
    \caption{Pre-emphasis and de-emphasis around the FM link.}
    \label{fig:pre_de_emphasis_chain}
\end{figure}

\textbf{Purpose:} \cref{fig:pre_de_emphasis_chain} shows that the boost and the cut sit on either side of the channel, so only the channel noise is filtered.

\subsection{Comparison}
\label{sec:emphasis_comparison}

\begin{tabularx}{\textwidth}{|c|y{2.8cm}|L|}
    \hline
    \rowcolor{gray!15}
    \cc{Stage}         &
    \yy{Filter}{2.8cm} &
    \CC{Effect}                                                                                     \\ \hline

    Tx                 &
    High-pass boost    &
    Raises the amplitude of high audio frequencies before the FM modulator.                         \\ \hline

    Rx                 &
    Low-pass cut       &
    Restores the original spectrum and attenuates the high-frequency noise added by the channel.    \\ \hline
\caption{Role of each emphasis filter.}
\label{tab:emphasis_filter_roles}
\end{tabularx}

\subsection{Numerical Check}
\label{sec:emphasis_numerical_check}

\textbf{Command:}
\begin{lstlisting}[language=Python, caption={Corner frequency of a 75 µs network.}]
import math
RC = 75e-6                  # seconds
f1 = 1 / (2 * math.pi * RC)
print(f"f1 = {f1:.0f} Hz")
\end{lstlisting}

\textbf{Output:}
\begin{lstlisting}[caption={Script output.}]
f1 = 2122 Hz
\end{lstlisting}

The result matches \cref{sec:emphasis_transfer_functions}. The file \verb|emphasis_check.py| holds the script.

\begin{repocard}{Course References}

\textbf{Notes Repository:}
\href{https://github.com/nurmdabd/Temp-3501}{\textit{Temp-3501/}}

\vspace{2mm}

\textbf{Textbook:}
B.~P.~Lathi, \textit{Modern Digital and Analog Communication Systems}, Ch.~5.

\end{repocard}
